%═══════════════════════════════════════════════════════════════════════════════
\chapter{El problema de aproximación de funciones}
\label{ch:cap1}
%═══════════════════════════════════════════════════════════════════════════════

%───────────────────────────────────────────────────────────────────────────────
\section{Un campo que no se puede medir en todas partes}
\label{sec:motivacion1}
%───────────────────────────────────────────────────────────────────────────────

\dificultad{1}{2}

Imagine un experimento sencillo: se quiere conocer la distribución de temperatura
en una placa metálica cuadrada. Se dispone de diez termopares soldados en
posiciones fijas $\bm{x}_1, \ldots, \bm{x}_{10}$, cada uno reportando una medición
$T_i = T(\bm{x}_i)$. La pregunta física es inmediata: ¿qué temperatura tiene
la placa en los miles de puntos donde no hay sensor?

Este problema aparece en formas distintas a lo largo de toda la física:
\begin{itemize}
  \item Un sismógrafo registra el movimiento del suelo en una estación; se
        quiere reconstruir el campo de ondas en la región circundante.
  \item Un detector registra la energía depositada en celdas discretas de un
        calorímetro; se quiere reconstruir el perfil continuo de la cascada.
  \item Un sistema de instrumentación mide la potencia neutrónica en unos
        pocos canales de un reactor; se quiere estimar el flujo $\phi(\bm{x})$
        en toda la geometría.
\end{itemize}

En todos estos casos el problema tiene la misma estructura matemática:
\emph{dado un número finito de evaluaciones de una función desconocida,
construir una función que las reproduzca y tenga comportamiento razonable
entre ellas}. Este es el \textbf{problema de aproximación de funciones},
y es el punto de partida del libro.

Lo que distingue al aprendizaje automático del análisis numérico clásico no
es la estructura del problema sino la \emph{escala}: en lugar de diez sensores
pueden ser diez millones de píxeles, y en lugar de una función escalar sobre
$\mathbb{R}^2$ puede ser un mapa de $\mathbb{R}^{10^6}$ en $\mathbb{R}^{10^3}$.
Las herramientas cambian por necesidad, pero el problema es el mismo.

\begin{fisicaconexion}{Aproximación vs. interpolación}
En física experimental se distingue entre \textbf{interpolar} (pasar exactamente
por los datos medidos) y \textbf{ajustar} (pasar cerca de los datos, absorbiendo
el ruido). La interpolación es apropiada cuando los datos son exactos (simulaciones
numéricas, tablas de constantes físicas). El ajuste es apropiado cuando los datos
tienen error de medición. Las redes neuronales operan en el régimen de ajuste,
pero con familias de funciones mucho más expresivas que los polinomios o las
series de Fourier que el físico suele emplear.
\end{fisicaconexion}

%───────────────────────────────────────────────────────────────────────────────
\section{Vectores, normas y el lenguaje del error}
\label{sec:algebra}
%───────────────────────────────────────────────────────────────────────────────

\dificultad{1}{1}

Para hablar de \emph{cuánto se equivoca} una aproximación necesitamos medir
distancias. En física esto es natural: la distancia entre dos vectores de posición,
la diferencia entre una solución numérica y la exacta, el residuo de una ecuación
algebraica. La herramienta matemática es la norma.

\begin{definicion}{Vectores y $p$-normas}
\label{def:normas}
Un \textbf{vector} es un elemento $\bm{x} = (x_1, \ldots, x_n)^\top \in
\mathbb{R}^n$. Dado $p \geq 1$, la \textbf{$p$-norma} de $\bm{x}$ es:
\[
  \|\bm{x}\|_p = \left(\sum_{i=1}^n |x_i|^p\right)^{1/p}.
\]
Tres casos de especial relevancia:
\[
  \|\bm{x}\|_1 = \sum_{i=1}^n |x_i|, \qquad
  \|\bm{x}\|_2 = \sqrt{\sum_{i=1}^n x_i^2}, \qquad
  \|\bm{x}\|_\infty = \max_{1\leq i\leq n} |x_i|.
\]
\end{definicion}

La elección de la norma no es cosmética: define qué tipo de error se penaliza.
La norma $2$ penaliza errores grandes de forma cuadrática (es sensible a
valores atípicos). La norma $1$ penaliza linealmente (más robusta). La norma
$\infty$ penaliza únicamente el peor error puntual.

\begin{fisicaconexion}{Las tres normas en física}
La norma $2$ aparece en la energía potencial elástica ($U = \frac{1}{2}k\|\bm{x}\|_2^2$)
y en la raíz del error cuadrático medio de instrumentación.
La norma $1$ es la distancia de Manhattan, relevante en transporte y en
estimadores robustos de señales con ruido impulsivo.
La norma $\infty$ es la norma del supremo: mide la desviación máxima, útil
en especificaciones de tolerancia en ingeniería.
\end{fisicaconexion}

\begin{definicion}{Producto interno euclídeo}
\label{def:inner}
Sean $\bm{x}, \bm{y} \in \mathbb{R}^n$. El \textbf{producto interno euclídeo} es:
\[
  \langle \bm{x}, \bm{y} \rangle = \bm{x}^\top \bm{y} = \sum_{i=1}^n x_i y_i,
\]
con $\|\bm{x}\|_2^2 = \langle \bm{x}, \bm{x} \rangle$.
La \textbf{desigualdad de Cauchy-Schwarz} establece:
\[
  |\langle \bm{x}, \bm{y} \rangle| \leq \|\bm{x}\|_2 \|\bm{y}\|_2,
\]
con igualdad si y solo si $\bm{x}$ e $\bm{y}$ son colineales.
El ángulo $\theta$ entre $\bm{x}$ e $\bm{y}$ se define por
$\cos\theta = \langle\bm{x},\bm{y}\rangle/(\|\bm{x}\|_2\|\bm{y}\|_2)$.
\end{definicion}

\begin{observacion}{Generalización funcional}
En espacios de funciones, el producto interno $\langle f, g\rangle =
\int_a^b f(x)g(x)\,dx$ es la generalización directa. La desigualdad de
Cauchy-Schwarz sigue valiendo y da lugar a la desigualdad de Schwarz en $L^2$.
Este producto interno es el fundamento del análisis espectral en mecánica cuántica
y del cálculo del gradiente en los capítulos siguientes.
\end{observacion}

%───────────────────────────────────────────────────────────────────────────────
\section{Familias paramétricas: el lenguaje unificador}
\label{sec:familias}
%───────────────────────────────────────────────────────────────────────────────

\dificultad{1}{2}

Un físico trabaja permanentemente con familias paramétricas sin llamarlas así.
La distribución de Maxwell-Boltzmann $f(v; T) = 4\pi n(m/2\pi k_B T)^{3/2}
v^2 e^{-mv^2/2k_BT}$ es una familia con parámetros $(n, T)$. La solución
de la ecuación de Schrödinger para el oscilador armónico depende de $\omega$.
La diferencia en aprendizaje automático es que los parámetros no tienen
interpretación física conocida a priori: se determinan optimizando un criterio
sobre los datos.

\begin{definicion}{Familia paramétrica de funciones}
\label{def:familia}
Una \textbf{familia paramétrica} es un mapa
$f: \mathbb{R}^n \times \mathcal{P} \to \mathbb{R}$
donde $\mathcal{P}$ es el espacio de parámetros. Para cada $\bm{\theta} \in
\mathcal{P}$ fijo, $f(\,\cdot\,;\bm{\theta}): \mathbb{R}^n \to \mathbb{R}$
es una función específica. La familia es la colección $\{f(\,\cdot\,;\bm{\theta})
: \bm{\theta} \in \mathcal{P}\}$.
\end{definicion}

\begin{ejemplo}{El oscilador amortiguado como familia paramétrica}
La posición de un oscilador amortiguado es
$x(t; A, \omega, \gamma) = A\,e^{-\gamma t}\cos(\omega t)$.
Aquí $t$ es la variable (lo que se observa), y $(A, \omega, \gamma)$ son los
parámetros (lo que se ajusta a los datos experimentales del oscilador).
En ML la estructura es idéntica: los parámetros no son frecuencias ni
amortiguamientos, sino pesos de una red, pero el lenguaje matemático es el mismo.
\end{ejemplo}

%───────────────────────────────────────────────────────────────────────────────
\section{Datos, pérdida y el problema de aproximación}
\label{sec:datos_perdida}
%───────────────────────────────────────────────────────────────────────────────

\dificultad{1}{2}

Con el lenguaje de las familias paramétricas, el problema de la sección de
apertura se formula con precisión.

\begin{definicion}{Conjunto de datos y pérdida empírica}
\label{def:datos_loss}
Un \textbf{conjunto de datos} es $\mathcal{D} = \{(\bm{x}_i, y_i)\}_{i=1}^N$
con $\bm{x}_i \in \mathbb{R}^n$ (entradas) e $y_i \in \mathbb{R}$ (salidas).
Dada una \textbf{función de pérdida} $\ell: \mathbb{R} \times \mathbb{R}
\to \mathbb{R}_{\geq 0}$ con $\ell(\hat{y}, y) = 0 \Leftrightarrow \hat{y} = y$,
la \textbf{pérdida empírica} es:
\[
  \mathcal{L}(\bm{\theta}) = \frac{1}{N}\sum_{i=1}^N \ell(f(\bm{x}_i;\bm{\theta}), y_i).
\]
Las dos pérdidas más comunes son el \textbf{error cuadrático medio}
$\ell_{\mathrm{MSE}}(\hat{y},y) = (\hat{y}-y)^2$ y el \textbf{error absoluto
medio} $\ell_{\mathrm{MAE}}(\hat{y},y) = |\hat{y}-y|$.
\end{definicion}

\begin{definicion}{Problema de aproximación}
\label{def:problema}
El \textbf{problema de aproximación} consiste en encontrar:
\[
  \bm{\theta}^* = \argmin_{\bm{\theta} \in \mathcal{P}}\,\mathcal{L}(\bm{\theta}).
\]
\end{definicion}

Esta definición unifica en un solo enunciado todos los problemas de ajuste
que el físico conoce. Lo que varía es la elección de la familia $f$, la
pérdida $\ell$ y el método para encontrar $\bm{\theta}^*$.

\begin{fisicaconexion}{Por qué MSE y no MAE}
El error cuadrático medio $\ell_{\mathrm{MSE}}$ penaliza los errores grandes
cuadráticamente. En términos estadísticos, minimizar la pérdida MSE equivale
a encontrar la \emph{media condicional} $\mathbb{E}[y|\bm{x}]$. Minimizar MAE
equivale a encontrar la \emph{mediana condicional}. En física experimental,
donde el ruido de medición es gaussiano (por el teorema central del límite),
MSE es el estimador de máxima verosimilitud. Cuando el ruido tiene colas
pesadas (ruido de Cauchy, interferencias impulsivas), MAE es más robusto.
El Capítulo~7 formaliza esta conexión entre pérdidas y modelos probabilísticos.
\end{fisicaconexion}

%───────────────────────────────────────────────────────────────────────────────
\section{Familias lineales en parámetros: solución exacta}
\label{sec:lineales}
%───────────────────────────────────────────────────────────────────────────────

\dificultad{2}{2}

Cuando la familia paramétrica tiene la forma $f(\bm{x};\bm{\theta}) =
\sum_{j=1}^p \theta_j \phi_j(\bm{x})$, los parámetros aparecen linealmente
y el problema de aproximación tiene solución cerrada.

\begin{definicion}{Familia lineal en parámetros y matriz de diseño}
\label{def:lineal}
Una familia es \textbf{lineal en los parámetros} si existen funciones fijas
$\phi_1, \ldots, \phi_p: \mathbb{R}^n \to \mathbb{R}$ (llamadas \textbf{funciones
base}) tales que $f(\bm{x};\bm{\theta}) = \bm{\theta}^\top \bm{\Phi}(\bm{x})$
con $\bm{\Phi}(\bm{x}) = (\phi_1(\bm{x}),\ldots,\phi_p(\bm{x}))^\top$.
La \textbf{matriz de diseño} $\bm{\Phi}_{\mathcal{D}} \in \mathbb{R}^{N\times p}$
tiene entradas $(\bm{\Phi}_{\mathcal{D}})_{ij} = \phi_j(\bm{x}_i)$.
La pérdida empírica MSE se escribe $\mathcal{L}(\bm{\theta}) = \frac{1}{N}
\|\bm{\Phi}_{\mathcal{D}}\bm{\theta} - \bm{y}\|_2^2$.
\end{definicion}

La tabla siguiente muestra las familias lineales más comunes en física
y sus matrices de diseño:

\begin{center}
\begin{tabular}{llll}
\toprule
\textbf{Base} & $\phi_j(\bm{x})$ & \textbf{Contexto físico} &
\textbf{Matriz $\bm{\Phi}_{\mathcal{D}}$} \\
\midrule
Monomios & $x^{j-1}$ & Ajuste de curvas & Vandermonde \\
Legendre & $P_{j-1}(x)$ & Expansión multipolar & Casi diagonal \\
Fourier & $e^{2\pi i (j-1)k/N}$ & Espectroscopia, señales & Unitaria \\
RBF gaussiana & $e^{-\|\bm{x}-\bm{c}_j\|^2/2\sigma^2}$ &
  Interpolación espacial, meshless & Densa \\
B-splines & $B_j(x)$ & Perfiles de flujo suaves & Banda \\
\bottomrule
\end{tabular}
\end{center}

\begin{proposicion}{Ecuaciones normales}
\label{prop:normales}
Si $\mathrm{rango}(\bm{\Phi}_{\mathcal{D}}) = p$, el minimizador de
$\mathcal{L}(\bm{\theta})$ es único y satisface:
\[
  \bm{\Phi}_{\mathcal{D}}^\top \bm{\Phi}_{\mathcal{D}}\,\bm{\theta}^*
  = \bm{\Phi}_{\mathcal{D}}^\top \bm{y},
  \qquad
  \bm{\theta}^* = (\bm{\Phi}_{\mathcal{D}}^\top \bm{\Phi}_{\mathcal{D}})^{-1}
  \bm{\Phi}_{\mathcal{D}}^\top \bm{y}
  =: \bm{\Phi}_{\mathcal{D}}^\dagger \bm{y}.
\]
El mapa $\bm{\Phi}_{\mathcal{D}}^\dagger$ se denomina
\textbf{pseudoinversa de Moore-Penrose}.
\end{proposicion}

\begin{proof}
$\mathcal{L}(\bm{\theta})$ es convexa y diferenciable. Su gradiente es
$\nabla_{\bm{\theta}}\mathcal{L} = \frac{2}{N}\bm{\Phi}_{\mathcal{D}}^\top
(\bm{\Phi}_{\mathcal{D}}\bm{\theta}-\bm{y})$.
Igualando a cero y usando que $\bm{\Phi}_{\mathcal{D}}^\top\bm{\Phi}_{\mathcal{D}}$
es invertible cuando el rango es $p$.
\end{proof}

\begin{fisicaconexion}{La pseudoinversa como estimador óptimo}
La solución $\bm{\theta}^* = \bm{\Phi}_{\mathcal{D}}^\dagger \bm{y}$ es el
\emph{estimador de mínima varianza no sesgado} (BLUE, Best Linear Unbiased
Estimator) cuando el ruido en las mediciones $y_i$ es aditivo, gaussiano e
independiente con varianza uniforme. Este es el teorema de Gauss-Markov, que
el Capítulo~7 derivará formalmente desde la perspectiva bayesiana.
\end{fisicaconexion}

\subsection{Caso particular: interpolación polinomial}
\label{subsec:interpolacion}

\dificultad{2}{1}

Cuando la familia base es $\phi_j(x) = x^{j-1}$ (monomios) y $N = p$,
la matriz de diseño es cuadrada y el problema de aproximación exige
interpolación exacta.

\begin{definicion}{Matriz de Vandermonde}
\label{def:vandermonde}
Dados $N$ nodos distintos $x_1 < \cdots < x_N$, la \textbf{matriz de Vandermonde}
$\bm{V} \in \mathbb{R}^{N \times N}$ tiene entradas $V_{ij} = x_i^{j-1}$.
\end{definicion}

\begin{proposicion}{Interpolación como caso particular del problema de aproximación}
\label{prop:interpolacion}
Con la familia monomio de grado $N-1$ y pérdida MSE, el minimizador
$\bm{\theta}^*$ satisface $\bm{V}\bm{\theta}^* = \bm{y}$ (interpolación exacta)
y es único porque $\det(\bm{V}) = \prod_{i<j}(x_j - x_i) \neq 0$.
\end{proposicion}

\begin{observacion}{Cuándo falla la interpolación polinomial}
La interpolación con monomios en nodos equiespaciados sufre el
\textbf{fenómeno de Runge}: para grado $\geq 10$, las oscilaciones en los
extremos del intervalo crecen sin control. El número de condicionamiento de
$\bm{V}$ crece exponencialmente con $N$. Esto no es un defecto de la
interpolación sino de la base: los nodos de Chebyshev o las bases de Legendre
eliminan el problema. Esta sensibilidad a la base es la primera motivación
para buscar familias \emph{adaptativas} que elijan sus propias funciones
base según los datos.
\end{observacion}

%───────────────────────────────────────────────────────────────────────────────
\section{Familias no lineales en parámetros}
\label{sec:nolineales}
%───────────────────────────────────────────────────────────────────────────────

\dificultad{2}{3}

Las familias lineales tienen una limitación estructural: las funciones base
$\phi_j$ son fijas antes de ver los datos. Para una distribución de temperatura
con estructura desconocida, no sabemos a priori qué bases elegir. La solución
natural es dejar que las bases se adapten a los datos.

\begin{definicion}{Familia no lineal en parámetros}
\label{def:nolineal}
Una familia es \textbf{no lineal en los parámetros} si no puede escribirse en
la forma de la Definición~\ref{def:lineal} para ninguna elección de funciones
base independientes de $\bm{\theta}$.
\end{definicion}

El ejemplo canónico son las \textbf{funciones de base radial con centros
aprendibles}: $f(x;\bm{\theta}) = \sum_{j=1}^m w_j\,e^{-(x-c_j)^2/2\sigma_j^2}$,
donde $\bm{\theta} = (w_j, c_j, \sigma_j)$. Los centros $c_j$ y anchos $\sigma_j$
aparecen dentro de la exponencial: la familia no es lineal en $\bm{\theta}$.

La consecuencia inmediata es que la pérdida empírica deja de ser convexa.
Esto implica múltiples mínimos locales, dependencia de la inicialización
y necesidad de métodos iterativos. El Capítulo~3 desarrolla esos métodos.
La pregunta que surge naturalmente es: ¿vale la pena esa complejidad adicional?

%───────────────────────────────────────────────────────────────────────────────
\section{El teorema de aproximación universal}
\label{sec:tau}
%───────────────────────────────────────────────────────────────────────────────

\dificultad{3}{3}

La respuesta a la pregunta anterior es el resultado más importante de este
capítulo. La intuición primero: una red neuronal de una capa es una suma de
funciones de la forma $\sigma(\bm{w}_j^\top \bm{x} + b_j)$, donde $\sigma$
es una función no lineal fija (la activación) y $(\bm{w}_j, b_j)$ son
parámetros aprendibles. Cada término es una ``función de umbral suavizado''
orientada en la dirección $\bm{w}_j$. La pregunta es: ¿cuántas de estas
funciones hacen falta para aproximar cualquier función continua?

\begin{teorema}{Aproximación universal \normalfont{(Cybenko 1989; Hornik et al.\ 1991)}}
\label{thm:universal}
Sea $\sigma: \mathbb{R} \to \mathbb{R}$ continua y no constante.
Sea $K \subset \mathbb{R}^n$ compacto y $g: K \to \mathbb{R}$ continua.
Para todo $\varepsilon > 0$ existe $M \in \mathbb{N}$ y parámetros
$\bm{w}_j \in \mathbb{R}^n$, $b_j, v_j \in \mathbb{R}$ tales que:
\[
  \sup_{\bm{x}\in K}\,\left|g(\bm{x}) - \sum_{j=1}^M v_j\,\sigma(\bm{w}_j^\top
  \bm{x} + b_j)\right| < \varepsilon.
\]
\end{teorema}

La intuición física antes de la demostración: $\sigma(\bm{w}^\top\bm{x}+b)$
es una función que transiciona de $0$ a $1$ (para $\sigma$ sigmoide) en el
hiperplano $\bm{w}^\top\bm{x} + b = 0$. La diferencia $\sigma(\bm{w}^\top
\bm{x}+b_1) - \sigma(\bm{w}^\top\bm{x}+b_2)$ aproxima una función ``meseta''
localizada entre dos hiperplanos paralelos. Sumando suficientes mesetas en
distintas orientaciones y posiciones se puede aproximar cualquier función
continua sobre un compacto.

\begin{proof}[Idea de la demostración]
Se demuestra que las combinaciones lineales de $\sigma(\bm{w}^\top\bm{x}+b)$
son densas en $\mathcal{C}(K)$ con la norma del supremo. La herramienta
es el teorema de Hahn-Banach: si el subespacio generado no fuera denso,
existiría un funcional lineal continuo no nulo que se anula sobre él.
Calculando ese funcional con la representación de Riesz-Markov y usando
que $\sigma$ es no constante se llega a contradicción.
La demostración completa puede consultarse en Cybenko~\cite{cybenko1989approximation}.
\end{proof}

\begin{observacion}{Qué garantiza y qué no garantiza el Teorema~\ref{thm:universal}}
El teorema garantiza \emph{existencia}: para cualquier función continua y
cualquier tolerancia $\varepsilon$, existe una red que la aproxima. \emph{No}
dice cuántas neuronas $M$ se necesitan (puede ser exponencial en la dimensión),
\emph{no} dice cómo encontrar los parámetros, y \emph{no} dice qué ocurre
fuera de $K$ ni cuando los datos son limitados. Esas tres preguntas abiertas
estructuran el resto del libro: cuántos parámetros (Capítulo~2), cómo
encontrarlos (Capítulos~3--6), y qué ocurre con datos limitados (Capítulo~2).
\end{observacion}

\begin{fisicaconexion}{La universalidad como principio físico}
El teorema de aproximación universal es el análogo para redes neuronales del
teorema de Fourier: cualquier función periódica puede representarse como
suma de senos y cosenos; cualquier función continua en un compacto puede
representarse como suma de neuronas sigmoides. La diferencia crucial es que
la base de Fourier es \emph{fija} (senos y cosenos no dependen de los datos),
mientras que las neuronas son \emph{adaptativas}: los vectores $\bm{w}_j$
se orientan según la estructura de los datos. Para un campo de temperatura
con simetría cilíndrica, la base de Fourier en coordenadas cilíndricas es
óptima; pero si la simetría no se conoce a priori, la red neuronal la descubre
durante el entrenamiento.
\end{fisicaconexion}

%───────────────────────────────────────────────────────────────────────────────
\section{Resumen del capítulo}
\label{sec:resumen1}
%───────────────────────────────────────────────────────────────────────────────

Este capítulo estableció un único hilo: el problema de aprendizaje supervisado
es el mismo problema de aproximación que el físico conoce de análisis numérico,
visto a través de un lente más general. La Definición~\ref{def:problema}
---encontrar $\bm{\theta}^*$ que minimiza la pérdida empírica--- unifica la
interpolación de Lagrange, el ajuste por mínimos cuadrados y el entrenamiento
de una red neuronal.

La distinción fundamental es entre familias \emph{lineales} en los parámetros
(bases fijas, solución cerrada por ecuaciones normales, pérdida convexa) y
familias \emph{no lineales} (bases adaptativas, solución iterativa, pérdida
no convexa con múltiples mínimos). El Teorema de Aproximación Universal
justifica el costo de la no linealidad: las redes neuronales pueden representar
cualquier función continua con precisión arbitraria.

El fenómeno de Runge muestra por qué las bases fijas clásicas fallan en alta
dimensión. La solución no es una base fija mejor sino una que se adapte
automáticamente a la estructura de los datos. Ese es el programa de los
capítulos siguientes.

%───────────────────────────────────────────────────────────────────────────────
\section*{Ejercicios resueltos}
\addcontentsline{toc}{section}{Ejercicios resueltos}
%───────────────────────────────────────────────────────────────────────────────

\begin{ejercicio}
\label{ej:vandermonde_resuelto}
\textbf{(Interpolación de tres puntos).}
Sean los nodos $x_1 = 0$, $x_2 = 1$, $x_3 = 2$ con valores $y_1 = 1$,
$y_2 = 3$, $y_3 = 2$. Encontrar el polinomio interpolante de grado 2.

\medskip\noindent\textit{Solución.}
La familia monomio con $d = 2$ tiene $\bm{\theta} = (\theta_0, \theta_1,
\theta_2)^\top$ y $f(x;\bm{\theta}) = \theta_0 + \theta_1 x + \theta_2 x^2$.
La matriz de Vandermonde es:
\[
  \bm{V} = \begin{pmatrix} 1 & 0 & 0 \\ 1 & 1 & 1 \\ 1 & 2 & 4 \end{pmatrix},
  \quad \bm{y} = \begin{pmatrix}1\\3\\2\end{pmatrix}.
\]
Resolviendo $\bm{V}\bm{\theta} = \bm{y}$ por sustitución hacia adelante:
fila 1 da $\theta_0 = 1$; fila 2 da $\theta_1 + \theta_2 = 2$; fila 3 da
$2\theta_1 + 4\theta_2 = 1$, de donde $\theta_2 = -3/2$ y $\theta_1 = 7/2$.
El polinomio interpolante es $f^*(x) = 1 + \frac{7}{2}x - \frac{3}{2}x^2$.

\medskip\noindent\textit{Verificación:}
$f^*(0) = 1$ \OK, $f^*(1) = 1 + 3.5 - 1.5 = 3$ \OK, $f^*(2) = 1 + 7 - 6 = 2$ \OK.
\end{ejercicio}

\begin{ejercicio}
\label{ej:mincuad_resuelto}
\textbf{(Mínimos cuadrados lineales).}
Datos $\mathcal{D} = \{(1,2),(2,3),(3,5),(4,4)\}$. Ajustar una recta
$f(x;\bm{\theta}) = \theta_0 + \theta_1 x$ con pérdida MSE.

\medskip\noindent\textit{Solución.}
La matriz de diseño y el vector de salidas son:
\[
  \bm{\Phi}_{\mathcal{D}} = \begin{pmatrix}1&1\\1&2\\1&3\\1&4\end{pmatrix},
  \qquad \bm{y} = \begin{pmatrix}2\\3\\5\\4\end{pmatrix}.
\]
Calculando:
\[
  \bm{\Phi}^\top\bm{\Phi} = \begin{pmatrix}4&10\\10&30\end{pmatrix},
  \qquad \bm{\Phi}^\top\bm{y} = \begin{pmatrix}14\\39\end{pmatrix}.
\]
Resolviendo las ecuaciones normales:
$\theta_1 = (4\cdot 39 - 10\cdot 14)/(4\cdot 30 - 100) = 16/20 = 0.8$,
$\theta_0 = (14 - 10\cdot 0.8)/4 = 1.5$.

La recta ajustada es $f^*(x) = 1.5 + 0.8x$ con pérdida empírica
$\mathcal{L}(\bm{\theta}^*) = 0.45$.

\medskip\noindent\textit{Interpretación física:} si los datos son mediciones
de posición con error gaussiano, esta recta es el estimador de máxima
verosimilitud de la velocidad ($\theta_1 = 0.8$ u/s) y posición inicial
($\theta_0 = 1.5$ u) de una partícula en movimiento uniforme.
\end{ejercicio}

%───────────────────────────────────────────────────────────────────────────────
\section*{Ejercicios propuestos}
\addcontentsline{toc}{section}{Ejercicios propuestos}
%───────────────────────────────────────────────────────────────────────────────

\begin{ejercicio}
\label{ej:normas}
\textbf{(Geometría de las normas).}
Sea $\bm{x} = (3,-4)^\top \in \mathbb{R}^2$.
\begin{enumerate}[label=(\alph*)]
  \item Calcule $\|\bm{x}\|_1$, $\|\bm{x}\|_2$ y $\|\bm{x}\|_\infty$.
  \item Demuestre que para cualquier $\bm{x} \in \mathbb{R}^n$ se cumple
        $\|\bm{x}\|_\infty \leq \|\bm{x}\|_2 \leq \|\bm{x}\|_1$.
  \item Un experimento produce residuos $\bm{r} = (0.1, 0.1, \ldots, 0.1)^\top
        \in \mathbb{R}^{100}$. Calcule las tres normas. ¿Cuál es más informativa
        sobre la calidad del ajuste?
\end{enumerate}
\end{ejercicio}

\begin{ejercicio}
\label{ej:vandermonde_condicionamiento}
\textbf{(Condicionamiento de Vandermonde).}
Sea la familia polinomial de grado $d = 10$ en $[0,1]$.
\begin{enumerate}[label=(\alph*)]
  \item Con nodos equiespaciados $x_i = i/10$, $i=0,\ldots,10$: calcule
        (o estime) el número de condicionamiento $\kappa(\bm{V}) =
        \|\bm{V}\|_2\|\bm{V}^{-1}\|_2$.
  \item Repita con nodos de Chebyshev $x_i = \frac{1}{2}\bigl(1 +
        \cos\frac{(2i+1)\pi}{22}\bigr)$, $i=0,\ldots,10$.
  \item ¿Por qué un número de condicionamiento alto implica que pequeños
        errores en $\bm{y}$ producen grandes errores en $\bm{\theta}^*$?
        Relacione con el ruido de medición en un experimento real.
\end{enumerate}
\end{ejercicio}

\begin{ejercicio}
\label{ej:rbf_vs_polinomio}
\textbf{(RBF vs.\ polinomial en señales de potencia).}
Sea $g(t) = e^{-0.1t}\cos(2\pi t/5)$ muestreada en $N = 8$ instantes
uniformes en $[0, 10]$ s, con ruido gaussiano de desviación $\sigma = 0.05$.
\begin{enumerate}[label=(\alph*)]
  \item Ajuste con la familia polinomial de grados $d = 2, 4, 7$.
        Para cada grado, calcule $\mathcal{L}(\bm{\theta}^*)$ en los datos
        de entrenamiento y en 100 puntos intermedios (error de generalización).
  \item Ajuste con familia RBF gaussiana con centros en los nodos de
        muestreo y $\sigma_{\mathrm{RBF}} = 1$ s. Compare con el ajuste
        polinomial de grado 4.
  \item ¿Para qué valor de $d$ el polinomio interpola exactamente los datos
        (incluido el ruido)? ¿Es eso deseable? Anticipe la conexión con el
        sobreajuste del Capítulo~2.
\end{enumerate}
\end{ejercicio}

\begin{ejercicio}
\label{ej:nolineal_noconvexo}
\textbf{(No convexidad en familias adaptativas).}
Sea $f(x; w, c) = w\,e^{-(x-c)^2}$ y el dato único $(x_1, y_1) = (0, 1)$.
\begin{enumerate}[label=(\alph*)]
  \item Escriba $\mathcal{L}(w,c)$ explícitamente y calcule las derivadas
        parciales $\partial\mathcal{L}/\partial w$ y $\partial\mathcal{L}/\partial c$.
  \item Muestre que el sistema $\nabla\mathcal{L} = \bm{0}$ no es lineal.
  \item Trace las curvas de nivel de $\mathcal{L}(w,c)$ para $w\in[-1,3]$,
        $c\in[-2,2]$. ¿Cuántos mínimos locales hay? ¿Depende la localización
        del mínimo de la inicialización?
\end{enumerate}
Este ejercicio ilustra por qué los métodos del Capítulo~3 son necesarios.
\end{ejercicio}

%───────────────────────────────────────────────────────────────────────────────
\section*{Ejercicios computacionales}
\addcontentsline{toc}{section}{Ejercicios computacionales}
%───────────────────────────────────────────────────────────────────────────────

\begin{ejerciciocomp}{Interpolación y mínimos cuadrados desde cero}
\textbf{Objetivo:} implementar los dos métodos clásicos del capítulo
y observar el efecto del condicionamiento.

\medskip\noindent\textbf{Datos:}
$N = 6$ mediciones de temperatura en posiciones $x = [0.0, 0.2, 0.4, 0.6,
0.8, 1.0]$ m con valores $T = [20.1, 22.3, 27.8, 35.1, 40.2, 38.9]$ °C.

\medskip\noindent\textbf{Algoritmo:}
\begin{enumerate}
  \item Construir la matriz de Vandermonde $\bm{V} \in \mathbb{R}^{6\times 6}$
        para la base monomio.
  \item Resolver $\bm{V}\bm{\theta} = \bm{T}$ por eliminación gaussiana
        (sin usar la función de resolución de sistemas del lenguaje).
  \item Evaluar el polinomio interpolante en 100 puntos uniformes en $[0,1]$.
  \item Construir la matriz de diseño RBF con centros $c_j = x_j$ y
        $\sigma = 0.2$ m. Resolver por ecuaciones normales.
  \item Graficar ambas interpolaciones y los datos originales.
\end{enumerate}

\medskip\noindent\textbf{Verificación:}
Para el polinomio interpolante: $f^*(0.5) \approx 32.1$ °C
(entre los valores en $x=0.4$ y $x=0.6$).
Para el ajuste RBF: la suma de cuadrados de residuos debe ser $< 0.01$ °C$^2$.

\medskip\noindent\textbf{Extensión:}
Añadir ruido gaussiano $\mathcal{N}(0, 0.5^2)$ a las mediciones y comparar
cómo cambia cada ajuste. ¿Cuál es más robusto al ruido?
\end{ejerciciocomp}

\begin{ejerciciocomp}{Paisaje de pérdida: lineal vs.\ no lineal}
\textbf{Objetivo:} visualizar la diferencia entre pérdida convexa
(familia lineal) y no convexa (familia no lineal).

\medskip\noindent\textbf{Datos:}
$N = 5$ puntos: $(x_i, y_i) = \{(0,1),(1,3),(2,2),(3,4),(4,1)\}$.

\medskip\noindent\textbf{Algoritmo:}
\begin{enumerate}
  \item Para la familia lineal $f(x;\bm{\theta}) = \theta_0 + \theta_1 x$:
        calcular $\mathcal{L}(\theta_0, \theta_1)$ en una cuadrícula
        $\theta_0 \in [-1,5]$, $\theta_1 \in [-2,3]$ con 100 puntos por eje.
        Graficar las curvas de nivel. Verificar que hay un único mínimo.
  \item Para la familia no lineal $f(x;w,c) = w\,e^{-(x-c)^2/2}$:
        calcular $\mathcal{L}(w,c)$ en la misma cuadrícula. Graficar las
        curvas de nivel. Contar cuántos mínimos locales se observan.
  \item Calcular $\bm{\theta}^*$ para la familia lineal usando las
        ecuaciones normales. Marcar el mínimo en el gráfico de nivel.
\end{enumerate}

\medskip\noindent\textbf{Verificación:}
Para la familia lineal: $\bm{\theta}^* \approx (1.4, 0.4)^\top$ con
$\mathcal{L}(\bm{\theta}^*) \approx 1.64$.

\medskip\noindent\textbf{Extensión:}
¿Qué ocurre con el paisaje de la familia no lineal si se aumenta $N$ a 20
puntos muestreados de la función $g(x) = 2\,e^{-(x-1.5)^2/2}$?
¿Desaparecen los mínimos locales?
\end{ejerciciocomp}

%───────────────────────────────────────────────────────────────────────────────
\section*{\texorpdfstring{$\bigstar$}{[*]}\ Estrella Polar: reconstrucci\'on del flujo neutr\'onico}
\addcontentsline{toc}{section}{\texorpdfstring{$\bigstar$}{[*]}\ Estrella Polar: flujo neutr\'onico}
\label{sec:estrella1}
%───────────────────────────────────────────────────────────────────────────────

\begin{observacion}{Nota al lector}
Esta sección es independiente del texto principal y puede omitirse sin
afectar la continuidad del libro. Aplica los conceptos del capítulo al
problema de reconstrucción del flujo neutrónico en un reactor de geometría
cúbica. No se asume conocimiento previo de física nuclear.
\end{observacion}

\bigskip

\begin{estrella}{Reconstrucción del flujo neutrónico desde detectores discretos}

\textbf{El sistema físico.}
Un reactor nuclear de investigación de geometría cúbica contiene un conjunto
de detectores de flujo neutrónico en posiciones fijas $\{\bm{x}_i\}_{i=1}^N
\subset [0,L]^3$. Cada detector reporta $\phi_i \approx \phi(\bm{x}_i)$,
donde $\phi(\bm{x})$ es el flujo neutrónico (neutrones cm$^{-2}$s$^{-1}$).
El objetivo es reconstruir $\phi(\bm{x})$ en toda la geometría del núcleo.

\textbf{Por qué importa.}
El flujo neutrónico determina la tasa de reacciones de fisión
$R(\bm{x}) = \Sigma_f(\bm{x})\phi(\bm{x})$ y por tanto la distribución
de potencia. Un operador necesita conocer $\phi(\bm{x})$ en toda la geometría
para verificar que la distribución de potencia no excede límites de seguridad
en ningún punto.

\textbf{El problema de aproximación.}
El flujo estacionario satisface la ecuación de difusión:
\[
  -\nabla\cdot D(\bm{x})\nabla\phi(\bm{x}) + \Sigma_a(\bm{x})\phi(\bm{x})
  = \frac{1}{\keff}\nu\Sigma_f(\bm{x})\phi(\bm{x}),
\]
con condiciones de contorno $\phi = 0$ en la frontera (aproximación de frontera
extrapolada). Esta EDP tiene solución analítica en geometría cúbica homogénea:
\[
  \phi^*(\bm{x}) = A\,\sin\!\left(\frac{\pi x}{L}\right)
  \sin\!\left(\frac{\pi y}{L}\right)\sin\!\left(\frac{\pi z}{L}\right).
\]
Esta solución es exactamente de la forma de la familia de funciones base de
la Definición~\ref{def:lineal}: es una sola función base fija con un coeficiente
$A$ a determinar. Cuando el reactor es heterogéneo (barras de control, reflector,
distintos enriquecimientos), la solución analítica no existe y hay que recurrir
a familias adaptativas. Ese es el problema que los capítulos siguientes resuelven.

\textbf{Conexión con el capítulo.}
La reconstrucción del flujo desde mediciones de detectores es el problema de
aproximación de la Definición~\ref{def:problema} con:
\begin{itemize}
  \item Familia: funciones base sinusoidales (para el reactor homogéneo) o
        RBF gaussianas (para el reactor heterogéneo).
  \item Datos: $\mathcal{D} = \{(\bm{x}_i, \phi_i)\}_{i=1}^N$ con $N$ típicamente
        entre 6 y 30 detectores.
  \item Pérdida: MSE con posible regularización para garantizar suavidad física.
\end{itemize}
El fenómeno de Runge tiene aquí una interpretación física directa: un polinomio
de grado alto interpola exactamente en los detectores pero oscila sin control
entre ellos, produciendo ``flujos negativos'' físicamente imposibles.
La regularización del Capítulo~2 es la herramienta que impone la suavidad
física como restricción matemática.
\end{estrella}

\begin{ejerciciocomp}{Reconstrucción del flujo en 1D}
\textbf{Datos sintéticos:}
Flujo teórico en un reactor cilíndrico de radio $R = 30$ cm:
$\phi_{\rm true}(r) = \phi_0 J_0(2.405 r/R)$
donde $J_0$ es la función de Bessel de orden cero y $\phi_0 = 10^{13}$
n cm$^{-2}$ s$^{-1}$.

Simular $N = 5$ mediciones en $r = [3, 9, 15, 21, 27]$ cm con ruido relativo
del 2\%: $\phi_i = \phi_{\rm true}(r_i)(1 + 0.02\,\varepsilon_i)$,
$\varepsilon_i \sim \mathcal{N}(0,1)$.

\medskip\noindent\textbf{Algoritmo:}
\begin{enumerate}
  \item Ajustar con familia polinomial de grados $d = 2, 4$.
  \item Ajustar con familia RBF con centros en los puntos de medición.
  \item Calcular el error máximo $\|\phi_{\rm aprox} - \phi_{\rm true}\|_\infty$
        en 100 puntos interiores para cada método.
  \item Verificar que el flujo reconstruido es no negativo en todo $[0, R]$.
\end{enumerate}

\medskip\noindent\textbf{Verificación:}
El máximo del flujo de Bessel ocurre en $r = 0$: $\phi_{\rm true}(0) = \phi_0$.
El mínimo es $\phi_{\rm true}(R) = 0$ (condición de contorno).

\medskip\noindent\textbf{Parámetros físicos de referencia:}
Para U-235, $\phi_0 \approx 10^{13}$ n cm$^{-2}$ s$^{-1}$ es un valor
típico para un reactor de investigación de 10 MW.
\end{ejerciciocomp}

%───────────────────────────────────────────────────────────────────────────────
\begin{notasbib}
El teorema de aproximación universal fue demostrado independientemente por
Cybenko~\cite{cybenko1989approximation} para funciones de activación sigmoide
y por Hornik et al.~\cite{hornik1989multilayer} para una clase más amplia de
activaciones. Ambas demostraciones usan el teorema de Hahn-Banach y la
representación de Riesz-Markov; la presentación de este capítulo sigue la
ruta de Cybenko. La extensión a redes profundas (múltiples capas), con tasas
de aproximación más eficientes que la red de una capa, fue desarrollada
posteriormente por Barron (1993) y Mhaskar (1996).

El fenómeno de Runge fue publicado por Carl Runge en 1901 usando la función
$f(x) = 1/(1+25x^2)$ como ejemplo. Los nodos de Chebyshev como solución
al problema de condicionamiento de Vandermonde aparecen en cualquier texto
estándar de análisis numérico; la referencia clásica es Trefethen (2019).

La pseudoinversa de Moore-Penrose y las ecuaciones normales son el fundamento
del método de mínimos cuadrados; su conexión con el estimador de Gauss-Markov
se desarrolla en el Capítulo~7 de este libro.

Para la ecuación de difusión neutrónica y la solución analítica en geometría
cúbica, la referencia estándar es Duderstadt \& Hamilton~\cite{duderstadt1976nuclear}.
Los métodos meshless basados en RBF para resolver la ecuación de difusión
neutrónica son un área activa de investigación;
véase Fasshauer~\cite{fasshauer2007meshfree} para el marco general.
\end{notasbib}
